\documentclass[10pt,a4paper]{amsart}
\usepackage[margin=19mm]{geometry}
\usepackage{amsmath,amssymb,mathtools}
\usepackage[scr=rsfs,cal=euler]{mathalfa}
\usepackage{xfrac}
\usepackage[unicode,hidelinks,bookmarks=false]{hyperref}
\newcommand{\ie}{i.e.\ }
\newcommand{\cf}{cf.\ }
\newcommand{\resp}{respectively\ }
\newcommand{\Subsection}{\subsection}
\providecommand{\nobreakdash}{\nobreak\hbox{-}}
\setlength{\emergencystretch}{2em}
\multlinegap0pt
\usepackage{bm}
\usepackage{bbm}
\usepackage{dsfont}
\usepackage{xspace}
\usepackage{cancel}
\usepackage{mathtools}
\usepackage[shortlabels]{enumitem}
\usepackage{amsthm}
\makeatletter
\@ifundefined{theorem}{\newtheorem{theorem}{Theorem}}{}
\@ifundefined{proposition}{\newtheorem{proposition}[theorem]{Proposition}}{}
\makeatother
\DeclareMathOperator{\cent}{cent}
\DeclareMathOperator{\diam}{diam}
\DeclareMathOperator{\rad}{rad}
\newcommand{\norm}[1]{\left|#1\right|}
\title[Chebyshev centers and radii for sets induced by quadratic ma]{Chebyshev centers and radii for sets induced by quadratic matrix inequalities}
\author{Amir Shakouri, Henk J. van Waarde, M. Kanat Camlibel}
\date{Source: arXiv:2403.05315v4 (CC BY 4.0)}
\begin{document}
\maketitle

\noindent\emph{Source and licence.} Chebyshev centers and radii for sets induced by quadratic matrix inequalities. Version arXiv:2403.05315v4 (CC BY 4.0), distributed under the Creative Commons Attribution 4.0 International licence, as recorded in the arXiv licence metadata for this exact version and shown on the abstract page https://arxiv.org/abs/2403.05315v4. Full rights notice: licenses/arxiv-2403.05315-CC-BY-4.0.txt.\par\smallskip

\noindent\textit{Editorial scope.} Counted unit: the Proposition whose source label is \texttt{prop:chebuniq} in the article (source0.tex lines 248--257, source label \texttt{prop:chebuniq}), delivered together with the article's own complete original proof of it (source0.tex lines 258--304, one \texttt{proof} environment of 47 lines). No mathematics is written, repaired or re-derived: statement and proof are the article's own text line for line. Required context retained uncounted: eq:diam (lines 230--233); eq:cheb (lines 235--238). Every definition, display and label of those lines is retained unchanged. Omitted from this package and not redistributed: 1--229, 234--234, 239--247, 305--1406. Nothing inside a retained excerpt is omitted or altered: every retained body is delivered exactly as the article's lines are written, so no omission receipt of any kind accompanies this package. Citations: the retained excerpts carry no citation command at all, so no bibliography entry of the article is transcribed and no reference is redistributed. Reference adaptation: none was needed and none was made; every reference inside the delivered unit resolves inside it, so no unresolved reference or citation is produced. Printed slips kept literal: no printed word, symbol, hypothesis or number of the counted statement or of its proof was altered. Layout and class: the article's own class is \texttt{sn-jnl} with packages including comment, caption, subcaption, enumitem, graphicx, multirow, amsmath, amssymb, amsfonts, mathtools, amsthm, mathrsfs, appendix, xcolor; the wrapper is the lane's pinned \texttt{amsart} editorial wrapper at 10pt on A4, which restates the theorem-like environments the retained text needs and adds the article's own macro definitions that the retained text needs (\texttt{\textbackslash cent}, \texttt{\textbackslash diam}, \texttt{\textbackslash norm}, \texttt{\textbackslash rad}), with extra wrapper packages (bm, bbm, dsfont, xspace, cancel, mathtools, enumitem). The delivered numbering is the unit's own. Assets: no retained span contains a figure environment, a graphics-inclusion command, a PDF-page inclusion, a table or a TikZ picture; no image file is copied into this package, the package declares no asset dependency and no image file is redistributed with it. Licence: version arXiv:2403.05315v4 of this article is distributed under the Creative Commons Attribution 4.0 International licence, as recorded in the arXiv licence metadata for this exact version and shown on the abstract page https://arxiv.org/abs/2403.05315v4. A full rights notice is delivered at licenses/arxiv-2403.05315-CC-BY-4.0.txt. Characters: every retained excerpt is pure ASCII, so the delivered engine, XeLaTeX with its default Latin Modern text font, reports no missing character.
\begin{equation}
\label{eq:diam}
\diam \mathcal{X}\coloneqq \max_{\Delta\in\mathcal{X}-\mathcal{X}} \norm{\Delta} 
\end{equation} 

\begin{equation}
\label{eq:cheb}
\rad \mathcal{X}\coloneqq \min_{C\in\mathbb{R}^{p\times q}} \max_{X\in\mathcal{X}} \norm{C-X}
\end{equation}

\begin{proposition}
\label{prop:chebuniq}
Let $\mathcal{X}\subset\mathbb{R}^{p\times q}$ be a compact set. Then:
\begin{enumerate}[label=\normalfont{(\alph*)},ref=\ref{prop:chebuniq}(\alph*)]
    \item\label{prop:chebuniq-a} $\tfrac{1}{2}\diam\mathcal{X}\leq\rad\mathcal{X}\leq\diam\mathcal{X}$.
    \item\label{prop:chebuniq-b} $\cent \mathcal{X}$ is compact and convex with no interior points.
    \item\label{prop:chebuniq-bc} For every $X_0\in\mathbb{R}^{p\times q}$ we have  $\diam(X_0+\mathcal{X})=\diam\mathcal{X}$, $\rad(X_0+\mathcal{X})=\rad\mathcal{X}$, and $\cent(X_0+\mathcal{X})=X_0+\cent\mathcal{X}$.
    \item\label{prop:chebuniq-c} If the norm is essentially strictly convex, then the set $\mathcal{X}$ has a unique Chebyshev center.
\end{enumerate}
\end{proposition}

\begin{proof}
(a) Observe that from the triangle inequality $\tfrac{1}{2}\norm{X_1-X_2}\leq \tfrac{1}{2}\norm{C-X_1}+\tfrac{1}{2}\norm{C-X_2}$. Therefore, taking $C$ to be a Chebyshev center, we have $\tfrac{1}{2}\norm{X_1-X_2}\leq \rad\mathcal{X}$ for all \mbox{$X_1,X_2\in\mathcal{X}$}, which implies $\tfrac{1}{2}\diam\mathcal{X}\leq\rad\mathcal{X}$. The second inequality is obvious from the definitions \eqref{eq:diam} and \eqref{eq:cheb}. 

(b) To prove the boundedness of $\cent\mathcal{X}$, we note that for any $C\in\cent\mathcal{X}$ and $X\in\mathcal{X}$ we have $\norm{C}\leq \norm{C-X}+\norm{X}\leq \rad\mathcal{X}+\norm{X}$. Since $\mathcal{X}$ is bounded, we see that $\cent\mathcal{X}$ is bounded. To prove closedness, consider a sequence $C_i\in\cent\mathcal{X}$, $i\in\mathbb{N}$, with $\lim_{i\rightarrow\infty} C_i=C$. Let $X\in\mathcal{X}$. Note that $\norm{C-X}\leq\norm{C-C_i}+\norm{C_i-X}\leq \norm{C-C_i}+\rad\mathcal{X}$. Take the limit as $i\rightarrow \infty$ to conclude that $\norm{C-X}\leq \rad\mathcal{X}$. Thus, $C\in\cent\mathcal{X}$ and therefore $\cent\mathcal{X}$ is closed. 
%Suppose that there exists $X\in\mathcal{X}$ such that $\norm{C-X}>\mathrm{rad}(\mathcal{X})$, which implies that there exists $\varepsilon>0$ such that $\norm{C-X}-\varepsilon>\mathrm{rad}(\mathcal{X})$. Since the sequence is convergent, there exists a $k\in\mathbb{N}$ such that $\norm{C_k-C}\leq\varepsilon$. Thus, for such a $k$, we have $\norm{C_k-X}\geq\norm{C-X}-\norm{C-C_k}\geq\norm{C-X}-\varepsilon>\mathrm{rad}(\mathcal{X})$, which contradicts with $C_k\in\mathrm{cent}(\mathcal{X})$ being a Chebyshev center. Therefore, the limit of any converging sequence in $\mathrm{cent}(\mathcal{X})$ belongs to $\mathrm{cent}(\mathcal{X})$, which implies its closeness. 
To prove convexity, suppose that $C_1,C_2\in\cent\mathcal{X}$. Then, by the triangle inequality, we have the following for any $\alpha\in[0,1]$:
\begin{equation}
\label{eq:prop:chebuniq-1}
\norm{\alpha C_1+(1-\alpha)C_2-X}\leq\alpha\norm{C_1-X}+(1-\alpha)\norm{C_2-X}\leq\rad\mathcal{X}.
\end{equation}
Therefore, for any $C_1,C_2\in\cent\mathcal{X}$ and any $\alpha\in[0,1]$ we have $\alpha C_1+(1-\alpha)C_2\in\cent\mathcal{X}$, which implies $\cent\mathcal{X}$ is convex. Now, we prove that $\cent\mathcal{X}$ has no interior points. Assume, on the contrary, that the interior of $\cent\mathcal{X}$ is nonempty. Let $C_1$ belong to the interior of $\cent\mathcal{X}$. Let $X\in\mathcal{X}$ be such that $|C_1-X|=\rad\mathcal{X}$. Since the interior of $\cent\mathcal{X}$ is nonempty, there exists a sufficiently small $\alpha>0$ such that $C_2=C_1+\alpha(C_1-X)\in\cent\mathcal{X}$. This implies that $|C_2-X|=(\alpha+1)|C_1-X|=(\alpha+1)\rad\mathcal{X}$. Since $\rad\mathcal{X}>0$, we have $|C_2 - X| > \rad \mathcal{X}$, meaning that $C_2$ is not a Chebyshev center, thus resulting in a contradiction. Therefore, $\cent\mathcal{X}$ does not have an interior point.

(c) For the diameter, the proof immediately follows from the fact that \mbox{$(X_0+\mathcal{X})-(X_0+\mathcal{X})=\mathcal{X}-\mathcal{X}$}. For the radius, we have
\begin{equation}
\rad (X_0+\mathcal{X})= \min_{C\in\mathbb{R}^{p\times q}} \max_{X\in\mathcal{X}} \norm{C-X-X_0}=\min_{C\in\mathbb{R}^{p\times q}+X_0} \max_{X\in\mathcal{X}} \norm{C-X},
\end{equation}
which is equal to $\rad \mathcal{X}$ as $\mathbb{R}^{p\times q}+X_0=\mathbb{R}^{p\times q}$. For the set of centers, we have 
\begin{equation}
\begin{split}
\cent(X_0+\mathcal{X})&=\{C:\norm{C-X}\leq\rad(X_0+\mathcal{X}) \ \text{for all}\  X\in X_0+\mathcal{X}\} \\
&= \{C:\norm{C-X-X_0}\leq\rad\mathcal{X} \ \text{for all}\  X\in\mathcal{X}\}\\
&=\{C+X_0:\norm{C-X}\leq\rad\mathcal{X} \ \text{for all}\   X\in\mathcal{X}\}=X_0+\cent\mathcal{X}.
\end{split}
\end{equation}

(d) Suppose that $C_1,C_2\in\cent\mathcal{X}$ and the norm $\norm{\ \cdot\ }$ is essentially strictly convex. We show that $C_1=C_2$. For this, recall from part (b) that $\cent\mathcal{X}$ is convex. Thus, we have $\tfrac{1}{2}(C_1+C_2)\in\cent\mathcal{X}$. Hence, there exists an $X\in\mathcal{X}$ such that
\begin{equation}
\label{eq:pf_equal_rad}
\norm{\tfrac{1}{2}(C_1+C_2)-X}=\rad\mathcal{X}.
\end{equation}
Moreover, we observe that $\norm{C_i-X}\leq\rad\mathcal{X}$ for $i=1,2$. Thus, it follows from the triangle inequality that
\begin{equation}
\label{eq:pf_inequal_rad}
\norm{\tfrac{1}{2}(C_1+C_2)-X}\leq\tfrac{1}{2}\norm{C_1-X}+\tfrac{1}{2}\norm{C_2-X}\leq\rad\mathcal{X}.
\end{equation}
From \eqref{eq:pf_equal_rad} and \eqref{eq:pf_inequal_rad} we have
\begin{equation}
\label{eq:prop:chebuniq-2}
\norm{\tfrac{1}{2}(C_1-X)+\tfrac{1}{2}(C_2-X)}=\tfrac{1}{2}\norm{C_1-X}+\tfrac{1}{2}\norm{C_2-X}=\rad\mathcal{X}.
\end{equation}
Since we have $\norm{C_i-X}\leq\rad\mathcal{X}$ for $i=1,2$, it follows from \eqref{eq:prop:chebuniq-2} that 
\begin{equation}
\label{eq:prop:chebuniq-3}
\norm{C_1-X}=\norm{C_2-X}=\rad\mathcal{X}.
\end{equation}
As the norm is essentially strictly convex, \eqref{eq:prop:chebuniq-2} implies that $C_1-X=\alpha(C_2-X)$ for some $\alpha\geq 0$. It follows now from \eqref{eq:prop:chebuniq-3} that $\alpha=1$. This implies that $C_1=C_2$.
\end{proof}

\end{document}
